\chapter{Ergebnisnachweise}
\label{app:pipelineergebnisse}
\begingroup
\raggedbottom
% Absatzüberschriften stehen im bestehenden Anhanggerüst direkt unter dem Kapitel.
\makeatletter\def\toclevel@subsection{1}\makeatother

% Interne Zuordnung der unveränderten Datenauszüge zu Originaldateien und
% Originalzeilen: data/pipelineergebnisse-f-2f/QUELLEN.json.
Dieser Anhang dokumentiert die Referenzrechnung der in
Abschnitt~\ref{sec:komponentenvalidierung} geprüften Lagerfrequenzen und ergänzt
Abschnitt~\ref{sec:pipelineergebnisse} um führende Parameterkonfigurationen,
Ergebnisse einzelner Lager und Messungen mit geänderten Fehlerindikationen im
DRS-Vergleich. Die Berechnung der DRS-Mittelungen ergänzt die Diskussion
in Abschnitt~\ref{sec:methodische_grenzen}.

\paragraphheading{Referenzrechnung der Lagerfrequenzen}
\label{app:lagerfrequenz_referenzrechnung}

% Eigene Nachrechnung anhand der Gleichungen~\eqref{eq:bpfo} und
% \eqref{eq:bpfi}. Eingabewerte: CWRU_6205_GEOMETRY und
% PADERBORN_MTK_IBU_6203_GEOMETRY aus config/bearings.py sowie die im
% Komponentenfall vorgegebenen Drehzahlen aus tests/test_peak_detection.py.
% Die Komponentenfunktion wurde für diese Nachrechnung nicht aufgerufen.
Die Vergleichswerte werden mit den Gleichungen~\eqref{eq:bpfo} und
\eqref{eq:bpfi} für die im Komponententest vorgegebenen Geometrien und
Drehzahlen berechnet.

Für den CWRU-Lagertyp 6205 beträgt die Drehfrequenz bei 1797~U/min
\[
  f_r=\frac{1797}{60}=29{,}95\,\mathrm{Hz}.
\]
Mit neun Wälzkörpern, einem
Wälzkörperdurchmesser von $7{,}94\,\mathrm{mm}$, einem Teilkreisdurchmesser
von $39{,}00\,\mathrm{mm}$ und einem Druckwinkel von $0^\circ$ ergibt sich:
\begin{equation}
  \begin{aligned}
    \mathrm{BPFO}
      &= \frac{9}{2}\cdot29{,}95\,\mathrm{Hz}
         \left(1-\frac{7{,}94}{39{,}00}\cos 0^\circ\right)
       \approx 107{,}3362\,\mathrm{Hz}
       \approx 107{,}34\,\mathrm{Hz}, \\
    \mathrm{BPFI}
      &= \frac{9}{2}\cdot29{,}95\,\mathrm{Hz}
         \left(1+\frac{7{,}94}{39{,}00}\cos 0^\circ\right)
       \approx 162{,}2138\,\mathrm{Hz}
       \approx 162{,}21\,\mathrm{Hz}.
  \end{aligned}
\end{equation}

Für den Paderborner Lagertyp MTK/IBU 6203 beträgt die Drehfrequenz bei
1500~U/min
\[
  f_r=\frac{1500}{60}=25{,}00\,\mathrm{Hz}.
\]
Mit acht Wälzkörpern, einem
Wälzkörperdurchmesser von $6{,}75\,\mathrm{mm}$, einem Teilkreisdurchmesser
von $29{,}05\,\mathrm{mm}$ und einem Druckwinkel von $0^\circ$ folgt:
\begin{equation}
  \begin{aligned}
    \mathrm{BPFO}
      &= \frac{8}{2}\cdot25{,}00\,\mathrm{Hz}
         \left(1-\frac{6{,}75}{29{,}05}\cos 0^\circ\right)
       \approx 76{,}7642\,\mathrm{Hz}
       \approx 76{,}76\,\mathrm{Hz}, \\
    \mathrm{BPFI}
      &= \frac{8}{2}\cdot25{,}00\,\mathrm{Hz}
         \left(1+\frac{6{,}75}{29{,}05}\cos 0^\circ\right)
       \approx 123{,}2358\,\mathrm{Hz}
       \approx 123{,}24\,\mathrm{Hz}.
  \end{aligned}
\end{equation}

\clearpage
\paragraphheading{Berechnung der DRS-Mittelungen}
\label{app:drs_mittelungen}

% Eigene, vom Autor beauftragte Nachrechnung der Implementierung:
% cm_pipeline_validation/core/drs.py, init_params, run_drs und estimate_filter.
% Eingaben: 64 kHz, vier Sekunden und Nenndrehzahlen 900/1500 U/min gemäß
% sec:datensatz_paderborn; Blockfaktor 5, Bezugsfrequenz 3000 Hz,
% Sicherheitsfaktor 3 und Periodenzahl 100 gemäß paderborn_summary.csv.
% Blockfaktor 20 ist ein Vergleichswert aus dem Suchraum, kein ausgewählter
% Evaluationsparameter. Auswertung der Indexfolge ohne Signalverarbeitung.
% Kontrolliert an N09_M07_F10_KI17_15.mat und N15_M07_F10_KI17_15.mat:
% tatsächliche Längen 256001/256002, Drehzahlen 900.0865868/1499.3208275,
% Blocklängen 21331/12805, jeweils ebenfalls 10/18 Paare bei Blockfaktor 5.
Für eine idealisierte Paderborner Aufzeichnung von vier Sekunden mit
$64\,\mathrm{kHz}$ gilt $L=256000$ Abtastwerte. Bei 1500~U/min
($25\,\mathrm{Hz}$) umfasst ein Block mit dem ausgewählten Faktor $B=5$
\[
  N=5\cdot\frac{64000}{25}=12800\ \text{Abtastwerte}
  \quad (0{,}2\,\mathrm{s}).
\]
Die freie Lücke zwischen den gepaarten Blöcken beträgt mit den festen
DRS-Parametern $\Delta=(3\cdot100/3000)\cdot64000=6400$
Abtastwerte, also $0{,}1\,\mathrm{s}$.

% Codebasis: run_drs, block_distance=N+delta und skip=block_distance;
% estimate_filter, range(skip, len(signal)-N, N), obere Grenze exklusiv.
% Aufeinanderfolgende Paare beginnen im Abstand N.
Der spätere Block jedes Paares beginnt laut Implementierung bei
$s_j=N+\Delta+jN<L-N$. Im Beispiel erfüllen $j=0,\ldots,17$
diese Bedingung: Es werden 18~Blockpaare gemittelt.
Tabelle~\ref{tab:drs_mittelungen} ergänzt die niedrigere Nenndrehzahl und
den Vergleich mit $B=20$. Blocklängen werden wie im Code auf ganze
Abtastwerte abgerundet.

% Eigene Nachrechnung derselben Indexfolge:
% B=5, 15 Hz: N=21333, delta=6400, m=10;
% B=20, 25 Hz: N=51200, delta=6400, m=3;
% B=20, 15 Hz: N=85333, delta=6400, m=1.
\begin{table}[H]
  \centering
  \caption{Rechnerische Zahl der DRS-Blockpaare bei vier Sekunden Messdauer
    und $64\,\mathrm{kHz}$; die freie Lücke beträgt jeweils $0{,}1\,\mathrm{s}$.}
  \label{tab:drs_mittelungen}
  \small
  \begin{tabular}{rrrrr}
    \hline
    \textbf{Blockfaktor} & \textbf{Drehzahl} & \textbf{Blocklänge} &
      \textbf{Blockdauer} & \textbf{Blockpaare} \\
    $B$ & (U/min) & (Abtastwerte) & (s) & \\
    \hline
    5 & 1500 & 12800 & 0,200 & 18 \\
    5 & 900 & 21333 & 0,333 & 10 \\
    20 & 1500 & 51200 & 0,800 & 3 \\
    20 & 900 & 85333 & 1,333 & 1 \\
    \hline
  \end{tabular}
\end{table}
Mit $B=20$ bleiben somit nur drei beziehungsweise ein Blockpaar übrig.
Die Rechnung bestimmt die Zahl der gemittelten Paare; ihre statistische
Unabhängigkeit und die erreichbare Diagnosegüte werden damit nicht geprüft.

\paragraphheading{Parameterkonfigurationen}

% Ergebnisbasis: paderborn_param_search.csv, alle 120 Zeilen; sämtliche
% summary-Dateien, Zeile 2; config/settings.py und utils/io.py,
% load_trained_parameters. Beide Tabellenlabels betreffen die realen
% Such- und Evaluationsparameter dieser Tabelle.
\begin{table}[H]
  \centering
  \caption{Untersuchter Suchraum und in beiden DRS-Varianten geladene Werte}
  \label{tab:ergebnisse_suchumfang}
  \label{tab:ergebnisse_vergleichsparameter}
  \small
  \begin{tabular}{llr}
    \hline
    \textbf{Parameter} & \textbf{Suchwerte} & \textbf{Evaluation} \\
    \hline
    DRS-Blockfaktor $B$ & 3, 5, 10, 15, 20 & 5 \\
    Schwellenfaktor $k$ & 1,5; 2; 3; 4 & 1,5 \\
    Relative Halbbreite $t$ (\%) & 1; 2 & 1 \\
    Lokale Halbbreite $w$ (Hz) & 10, 20, 40 & 20 \\
    \hline
    \multicolumn{3}{l}{Fest: Ordnungen 1 und 2, beide Peaks erforderlich;} \\
    \multicolumn{3}{l}{Zerlegungstiefe 8, DRS-Bezugsfrequenz 3000\,Hz,} \\
    \multicolumn{3}{l}{Sicherheitsfaktor 3 und Periodenzahl 100.} \\
    \hline
  \end{tabular}
\end{table}

% Ergebnisbasis: paderborn_param_search.csv, Originalzeilen 2, 3, 4 und 12.
% Auswahlzweck: drei führende Ränge; Rang 11 als Maximum von detection_f1.
Tabelle~\ref{tab:ergebnisse_rangliste} zeigt die drei führenden
Konfigurationen und Rang elf mit dem höchsten binären F1-Maß.
Alle Trainingskennzahlen beziehen sich auf dieselben 120 Messungen.

\begin{table}[H]
  \centering
  \caption{Ausgewählte Parameterkonfigurationen und Trainingskennzahlen}
  \label{tab:ergebnisse_rangliste}
  \small
  \setlength{\tabcolsep}{4pt}
  \begin{tabular}{rrrrrrrr}
    \hline
    Rang & $B$ & $k$ & $t$ (\%) & $w$ (Hz) &
      \shortstack{exakt\\(\%)} & \shortstack{F1 Diagnose\\(\%)} &
      \shortstack{F1 binär\\(\%)} \\
    \hline
    1 & 5 & 1,5 & 1,00 & 20 & 45,83 & 45,83 & 40,38 \\
    2 & 3 & 1,5 & 1,00 & 10 & 45,00 & 44,81 & 40,74 \\
    3 & 5 & 1,5 & 2,00 & 20 & 44,17 & 45,45 & 48,70 \\
    11 & 3 & 1,5 & 2,00 & 10 & 42,50 & 42,98 & 55,12 \\
    \hline
  \end{tabular}
\end{table}

\paragraphheading{Paderborner Lager}

% Ergebnisbasis: alle 120 Zeilen von paderborn_results.csv; Gruppierung
% nach bearing, expected und vollständigem detected.
\begin{table}[H]
  \centering
  \caption{Paderborner Ausgabemengen mit DRS, jeweils 20 Messungen pro Lager}
  \label{tab:ergebnisse_paderborn_lager}
  \small
  \begin{tabular}{llrrrr}
    \hline
    \textbf{Lager} & \textbf{Referenz} & normal & BPFO & BPFI & BPFO+BPFI \\
    \hline
    K003 & normal & 18 & 0 & 2 & 0 \\
    K004 & normal & 18 & 2 & 0 & 0 \\
    KA04 & BPFO & 0 & 19 & 0 & 1 \\
    KA16 & BPFO & 0 & 20 & 0 & 0 \\
    KI16 & BPFI & 0 & 0 & 11 & 9 \\
    KI17 & BPFI & 14 & 1 & 5 & 0 \\
    \hline
  \end{tabular}
\end{table}

\clearpage
\paragraphheading{DRS-Vergleich}

% Ergebnisbasis: paarweise Zuordnung über file in paderborn_results*.csv
% und cwru_paderborn_params_results*.csv; alle 13 geänderten Klassenmengen.
% P1--P8 und C1--C5 ordnen die ausgewählten Originalzeilen eindeutig zu.
Tabelle~\ref{tab:ergebnisse_einzelfaelle} enthält sämtliche Messungen,
deren Ausgabemenge sich zwischen den Varianten ändert. P bezeichnet
Paderborn, C den CWRU-Datensatz. Die übrigen 112 Paderborner und 35
CWRU-Ausgaben bleiben gleich. Die Auswahl erfasst sowohl Gewinne als auch
Verluste exakter Zuordnungen und zusätzliche Indikationen bei weiterhin
binär erkannten Schäden oder Fehlalarmen.

\begin{table}[H]
  \centering
  \caption{Geänderte Ausgabemengen bei gleicher Messung und Konfiguration}
  \label{tab:ergebnisse_einzelfaelle}
  \small
  \setlength{\tabcolsep}{3pt}
  \begin{tabular}{lllll}
    \hline
    \textbf{Fall} & \textbf{Messungskennung} & \textbf{Referenz} &
      \textbf{ohne DRS} & \textbf{mit DRS} \\
    \hline
    P1 & \path{N15_M07_F10_K003_1.mat} & normal & BPFO & normal \\
    P2 & \path{N15_M07_F04_K004_10.mat} & normal & normal & BPFO \\
    P3 & \path{N15_M07_F04_KA04_15.mat} & BPFO & BPFO+BPFI & BPFO \\
    P4 & \path{N15_M07_F04_KA04_20.mat} & BPFO & BPFO+BPFI & BPFO \\
    P5 & \path{N09_M07_F10_KI16_10.mat} & BPFI & BPFO+BPFI & BPFI \\
    P6 & \path{N09_M07_F10_KI16_15.mat} & BPFI & BPFI & BPFO+BPFI \\
    P7 & \path{N15_M01_F10_KI16_5.mat} & BPFI & BPFI & BPFO+BPFI \\
    P8 & \path{N09_M07_F10_KI17_15.mat} & BPFI & BPFI & normal \\
    \hline
    C1 & \path{IR021_1.mat} & BPFI & BPFO+BPFI & BPFI \\
    C2 & \path{normal_1.mat} & normal & BPFO & BPFO+BPFI \\
    C3 & \path{normal_2.mat} & normal & BPFO & BPFO+BPFI \\
    C4 & \path{OR007@12_1.mat} & BPFO & normal & BPFO \\
    C5 & \path{OR007@12_2.mat} & BPFO & BPFO & BPFO+BPFI \\
    \hline
  \end{tabular}
\end{table}

% Eigene Auswahl aus den Originalzeilen: P1 und C1 zeigen geänderte Ausgaben
% bei gleichem Band; P8 und C4 geänderte Ausgaben bei geändertem Band;
% P9 zeigt geändertes Band bei unveränderter Ausgabe. Frequenzen in Hz,
% Originalwerte in den CSV-Auszügen ungerundet erhalten.
Tabelle~\ref{tab:ergebnisse_baender} ergänzt ausgewählte Bandangaben.
P1 und C1 zeigen eine geänderte Ausgabe bei gleichem Band, P8 und C4 eine
geänderte Ausgabe bei verändertem Band. Als Gegenbeispiel dient P9,
\path{N15_M01_F10_K003_10.mat}: Diese fehlerfreie Referenzmessung bleibt
in beiden Varianten ohne Ringindikation, obwohl sich das gewählte Band
ändert. Ein Bandwechsel und eine geänderte Klassenmenge fallen somit
nicht durchgehend zusammen.

\begin{table}[H]
  \centering
  \caption{Ausgewählte Mittenfrequenzen $f_c$ und Bandbreiten $b$ in Hz}
  \label{tab:ergebnisse_baender}
  \small
  \begin{tabular}{lrrrr}
    \hline
    & \multicolumn{2}{c}{\textbf{ohne DRS}} &
      \multicolumn{2}{c}{\textbf{mit DRS}} \\
    \textbf{Fall} & $f_c$ & $b$ & $f_c$ & $b$ \\
    \hline
    P1 & 8666,67 & 1333,33 & 8666,67 & 1333,33 \\
    P8 & 750,00 & 500,00 & 500,00 & 1000,00 \\
    P9 & 8333,33 & 666,67 & 500,00 & 1000,00 \\
    C1 & 4500,00 & 1000,00 & 4500,00 & 1000,00 \\
    C4 & 2812,50 & 375,00 & 4500,00 & 1000,00 \\
    \hline
  \end{tabular}
\end{table}
\par\clearpage
\endgroup
